\documentclass{article}

% PACKAGES =====================================================
% font settings
\usepackage{titlesec}
\usepackage{lmodern}
\usepackage[T1]{fontenc} % For improved sans-serif font
% Change title format to sans-serif
\titleformat{\title}{\sffamily\LARGE\bfseries}{}{0pt}{}
% Change section format to sans-serif
\titleformat{\section}{\sffamily\Large\bfseries}{}{0pt}{}
% Change subsection format to sans-serif
\titleformat{\subsection}{\sffamily\large\bfseries}{}{0pt}{}
% Change subsubsection format to sans-serif
\titleformat{\subsubsection}{\sffamily\normalsize\bfseries}{}{0pt}{}
% page settings
\usepackage[a4paper,scale=0.8]{geometry}
% language settings
\usepackage[american]{babel}
% captions settings
\usepackage{caption}
\captionsetup{labelfont={small,sf,bf}}
\setlength\belowcaptionskip{5pt}
% graphics settings
\usepackage{graphicx}
% tweaking the footnotes
\usepackage[hang,flushmargin]{footmisc}
% bibliogrpahy
\usepackage[backend = biber, autolang=other, style = nature]{biblatex}
%\addbibresource{my_bib.bib}
% graph creation
\usepackage{tikz}
\usepackage{float}
% hyperref
\usepackage{hyperref}
% table settings
\usepackage{multirow}
\usepackage{csvsimple}
% float barrier
\usepackage{placeins}
% list settings
\usepackage{enumitem}
% figure settings
\usepackage{subcaption}

% code settings
\usepackage{listings}
\usepackage{xcolor}

\lstdefinestyle{mystyle}{
    backgroundcolor=\color{white},   % choose the background color
    basicstyle=\ttfamily\footnotesize, % size of fonts used for the code
    keywordstyle=\color{blue},	     % color of keywords
    commentstyle=\color{green!60!black}, % color of comments
    stringstyle=\color{red},	     % color of strings
    numberstyle=\tiny\color{gray},   % style of line numbers
    breaklines=true,		     % automatic line breaking
    captionpos=b, 	     % position of the caption
    frame=single, 	     % adds a frame around the code
    numbers=left, 	     % where to put the line numbers
    numbersep=5pt,		     % how far the line numbers are from the code
    showspaces=false,		     % show spaces adding particular underscores
    showstringspaces=false,	     % underline spaces within strings
    showtabs=false,		     % show tabs within strings adding particular underscores
    tabsize=2			     % sets default tabsize to 2 spaces
}
\lstset{style=mystyle}

% header settings
\usepackage{fancyhdr}
\pagestyle{fancy}
\fancyhf{} % Clear all header and footer fields
\fancyfoot[C]{\thepage} % Add page numbering to the center footer

\renewcommand{\headrulewidth}{0pt} % Remove horizontal line from header
% first page customization
\fancypagestyle{plain}{
\fancyhead[L]{\textit{Lisa PARUIT}- week3 assigment}
% Left header
%  \fancyhead[R]{
%    \raisebox{-0.3\height}{\includegraphics[height=1cm]{Premier_logo.jpg}}
%    \raisebox{-0.3\height}{\includegraphics[height=1cm]{Second_logo.jpg}}
%  }
}
% title settings
\title{\sffamily\LARGE\bfseries Is Florida getting warmer?}
\author{\sffamily Lisa PARUIT\\CMEECourseWork}
\date{}


% BEGIN DOC ====================================================

\begin{document}
\maketitle

% Changement du titre de l'environnement abstract
%\renewenvironment{abstract}
%{\begin{center}\textbf{Résumé}\end{center}}{}

%\begin{abstract}
%    Ecrire le résumé ici
%\end{abstract}

% La créationd 'un second environement abstract est possible en réitérant cette série de lignes

% =======================================================================

We have a time record of the average temperature in Florida over the years. This time series is ploted on \autoref{Fig1} where one can observe there seems to be a trend of rising temperature in Florida between. However, we want to make sure that the positive correlation between time and temperature is significant. We will therefore perform a simple statistical test, using correlation of randomized value of year and temperature as \textit{null hypothesis}. If a significant amount of random correlations are above the effective correlation between time and temperature, then it means that the difference with the null hypothesis is not significant. If the null hypothesis correlation set is large enough, the ratio of random correlations that are above the effective one is refered to as \textit{p-value}.

\begin{figure}[!h]
    \centering
    \includegraphics[width = 0.65\textwidth]{"../results/Florida_temperature_plot.png"}
    \caption{Time series of Florida temperature with linear regression line in red.}\label{Fig1}
\end{figure}

We therefore find out that there is a significant (p value is close to 0) positive correlation (0.533) between time and temperature, suggesting that Florida is indeed getting warmer.



%\printbibliography
\end{document}